\part{Ejecución: Aseguramiento de Calidad y Pulido Final}

\chapter{Introducción Técnica a la Validación de Componentes}

La estrategia de validación final adopta un enfoque orientado al aseguramiento de la calidad mediante pruebas de integración automáticas. A nivel de infraestructura, se utiliza el motor \textit{Jest} para instanciar el \textit{Game Service} y los servicios de usuarios junto con bases de datos aisladas en memoria. Esto permite simular con precisión el flujo asimétrico de la partida, certificando la invariabilidad algorítmica de las condiciones de victoria y mitigando la aparición de regresiones indeseadas tras refactorizaciones de última hora.

\textbf{Estructura de Directorios} \\
\begin{verbatim}
packages/
+-- server/
    +-- tests/
        |-- gameService.test.mjs
        |-- userRoutes.test.mjs
        +-- userService.test.mjs
\end{verbatim}

\textbf{Snippets Estratégicos de Código} \\
\begin{lstlisting}[language=Java,breaklines=true]
// packages/server/tests/gameService.test.mjs
describe('gameService', () => {
  test('startGame inicializa la partida correctamente', async () => {
    // Aserciones internas del motor logic
  });
});
\end{lstlisting}

\chapter{Iteración 11: Sistema de Acciones Pendientes, Interfaces Asíncronas y Modales de Cartas}

\section{Definición de Características}
El objetivo principal de este ciclo consiste en desarrollar la capa de interacción visual que soporta las mecánicas avanzadas del motor de juego. Se busca implementar las interfaces reactivas necesarias para gestionar las decisiones de los jugadores, permitiendo que el sistema de juego sea capaz de solicitar información adicional de forma dinámica ante el uso de determinadas cartas. Asimismo, se pretende integrar un sistema de control de la interfaz que coordine la visualización de elementos interactivos con el estado global de la partida, asegurando que el flujo de juego se mantenga coherente y que las acciones pendientes se presenten de forma intuitiva para el usuario. Esta etapa busca consolidar la experiencia de juego en el cliente, proporcionando las herramientas visuales necesarias para manejar la complejidad asimétrica de las interacciones.

\section{Diseño de la Solución}
Para cumplir el requisito de rendimiento (\textbf{RNF-REN}), el diseño plantea traducir la lógica del servidor a una experiencia reactiva. Se proyecta un diccionario visual unificado (\texttt{pendingActionsUI.tsx}) capaz de instanciar modales específicos (como el ataque \textit{Strike} del Rey) congelando la interfaz general sin perder el contexto visual del tablero.

\section{Detalles de Implementación}
El diccionario dinámico se implementa mapeando las claves de acción del servidor con sus respectivos componentes en \textit{React}. Cuando \textit{Socket.io} notifica un estado con \texttt{pendingAction}, el cliente evalúa si el foco recae en el jugador local. En caso afirmativo, se invoca el componente del diccionario forzando la decisión del usuario y bloqueando algorítmicamente el resto de interacciones (\textit{force play}) hasta que el controlador del servidor resuelva la promesa asíncrona.

Se implementó de manera robusta la sincronización de estas interfaces. A continuación, se muestra cómo se modela el comportamiento de la carta \texttt{THUG} tanto en el cliente para recolectar el objetivo, como en el servidor para resolver su lógica:

\begin{lstlisting}[basicstyle=\scriptsize\ttfamily,caption={Validación de UI para la carta THUG en el Frontend},label={cod:pendingActionsUI1},breaklines=true]
"THUG": {
    instructionText: "Selecciona 1 guardia en el pueblo rival para descartarlo",
    allowedZones: ['rivalTown'],
    canConfirm: (selectedCards) => {
        const validate = selectedCards.every(c => c.typeKing === 'King' && c.position === 'rivalTown');
        return selectedCards.length === 1 && validate;
    }, 
    formatPayload: (selectedCards) => {
        return { targetUid: selectedCards[0]?.uid || "" }
    },
}
\end{lstlisting}

\begin{lstlisting}[basicstyle=\scriptsize\ttfamily,caption={Resolución del efecto de THUG en el Backend},label={cod:rebelPendingActions1},breaklines=true]
"THUG": (gameState, targetData) => {
    const { targetUid } = targetData; 
    const cardIndex = gameState.players.king.town.findIndex(card => card.uid === targetUid);
    if (cardIndex === -1) throw new Error('Carta no encontrada en el pueblo');
    const [targetCard] = gameState.players.king.town.splice(cardIndex, 1);
    targetCard.isRevealed = true;
    gameState.discardPile.push(targetCard);
    drawCardFromDeck(gameState, 'peasant');
    changeTurn(gameState);
    return gameState;
}
\end{lstlisting}

\section{Seguimiento, Control y Replanificación}
Durante el acoplamiento del diccionario, el control de la interfaz detecta desajustes visuales (\textit{frontend bugs}) al intentar interactuar con el mazo general mientras un modal estaba activo. Este defecto fuerza una replanificación en la jerarquía de dependencias de \textit{React}, reescribiendo los \textit{hooks} para forzar el paso de turno únicamente cuando se confirma el retorno HTTP/WS 200 desde el servidor, estabilizando así el flujo. El diseño mediante diccionarios demostró ser muy escalable. Aplicando criterio técnico, se estandarizó este patrón para delegar su uso a todo el catálogo de cartas, marcando el objetivo de la siguiente iteración como un proceso de control de calidad sobre cada carta individual. La reescritura de los 	extit{hooks} demandó esfuerzo fuera de la estimación inicial, logrando mitigar la desviación. Se da por completada la tarea 4.3 (Sistema de interrupciones y acciones asíncronas) definida en la EDT.


\chapter{Iteración 12: Consolidación Funcional del Motor de Juego y Lógica de Cartas}

\section{Definición de Características}
El propósito de este ciclo consiste en completar y consolidar el conjunto de mecánicas operativas que componen el sistema de juego asimétrico. Se busca integrar la totalidad de las funcionalidades de las cartas de ambos bandos, garantizando que los efectos complejos y las acciones encadenadas se resuelvan de forma íntegra y sin ambigüedades. Asimismo, se pretende asegurar que el motor de juego sea capaz de gestionar mutaciones múltiples sobre el estado de la partida de forma atómica, manteniendo la coherencia de la información entre el cliente y el servidor en todo momento. Esta etapa busca finalizar la implementación de las reglas lógicas del juego, proporcionando una base funcional robusta que soporte la diversidad de estrategias disponibles para los jugadores.

\section{Diseño de la Solución}
Para abordar el volumen de trabajo de forma eficiente, se opta por un enfoque metodológico riguroso basado en una matriz de control de calidad. Se diseña un documento de seguimiento donde se listan pormenorizadamente todas las cartas que componen el juego junto a cada una de sus acciones y efectos posibles. Este diseño permite actuar el desarrollo no como una tarea abierta, sino como un proceso sistemático de verificación y cierre de requisitos funcionales mediante una activación, prueba y corrección procedimental.

\section{Detalles de Implementación}
La programación de acciones múltiples (ej. una carta que permite mirar la mano rival y descartar una carta) se implementa bifurcando el \textit{GameService}. El servidor valida la orden y crea el \texttt{pendingAction}, inyectando al cliente los datos temporales. Una vez tomada la decisión, la resolución final unifica las propiedades alteradas y reanuda el turno general purgando la caché diferida de \textit{Redis}, garantizando una ejecución inmutable.

La fase de ejecución técnica se desarrolla siguiendo un ciclo iterativo de \textit{Test-Driven Development} manual. Utilizando la lista definida en la fase de diseño, se procede a verificar cada acción individualmente en el entorno local, revisando los cambios tanto en Redis como en el cliente. El procedimiento se ejecuta bajo los siguientes pasos:
\begin{itemize}
    \item Recorrido sistemático de la lista probando cada acción y anotando su correcto funcionamiento o sus fallos.
    \item Por cada acción que presentaba anomalías, se procede a la modificación del código para corregir o terminar de implementar la funcionalidad restante.
    \item Re-ejecución de la prueba tras cada cambio en el código hasta validar la resolución de la acción, marcándola entonces como resuelta.
    \item Monitorización continua durante la partida para detector nuevos fallos en acciones previamente validadas, re-incorporándolas a la lista de revisión en caso necesario.
\end{itemize}

\section{Seguimiento, Control y Replanificación}
El control de calidad evidencia fricciones temporales al enlazar el diccionario visual con resoluciones múltiples encadenadas. Las respuestas de validación fallan. Como medida correctiva, se replanifica la arquitectura de respuestas del \textit{GameService} para homogeneizar los formatos de transferencia (DTO) en cada paso de validación intermedia, dejando el sistema de cartas prácticamente completado. La aplicación de este procedimiento resultó en un éxito técnico notable, logrando dejar operativa la gran mayoría de las cartas en la partida. La metodología sistemática permitió, además, identificar y catalogar errores específicos o funcionalidades pendientes de gran complejidad —como la mecánica de infiltración profunda— que no pudieron solventarse en este ciclo. Como medida de replanificación, estas necesidades aisladas se definieron como requerimientos independientes para ser abordados en las siguientes etapas operativas. Debido a esta fricción técnica, el avance funcional resulta parcial. Se aplaza la finalización de la tarea 4.2 de la EDT a la siguiente iteración, reprogramando la implementación del módulo de infiltración.


\chapter{Iteración 13: Estabilización del Bucle de Partida y Mecánica de Infiltración}

\section{Definición de Características}
El objetivo establecido para este ciclo consiste en estabilizar el ciclo de vida completo de las partidas, desarrollar las interfaces específicas para las mecánicas de infiltración y añadir la automatización del robo de cartas (Bando del Rey). Se busca asegurar que el bucle de juego sea capaz de gestionar correctamente las transiciones críticas entre las distintas fases o Eras, garantizando que el estado del tablero se reinicie y se asigne de forma determinista. Asimismo, se pretende implementar los componentes visuales necesarios para permitir que el jugador interactúe con el bando rival de forma controlada, proporcionando una experiencia fluida en el manejo de la información efímera. Esta etapa busca dotar al sistema de la robustez necesaria para soportar partidas completas, eliminando cualquier inconsistencia en la progresión lógica del juego.

\section{Diseño de la Solución}
Para asegurar la fiabilidad (\textbf{RNF-FIA}), el diseño requiere crear un componente reactivo (\texttt{InfiltrateModal.tsx}) acoplado al diccionario de acciones, dotando al jugador de una vista controlada para examinar la mano rival. Paralelamente, se diseña la estabilización de los scripts de entorno local (\texttt{npm run start:local}) y del bloque de aserciones de victoria (\textit{Win Condition}) en el servidor.

Para el robo automático, se diseña una lógica de interceptación en la transición de turnos. En lugar de una simple alternancia de roles, se proyecta una función de envoltura que evalúa el estado saliente: si el jugador que finaliza es el Rey, el sistema debe invocar la lógica de robo antes de actualizar el turno en Redis.

Respecto a la infiltración, se analizaron diversas alternativas de diseño para que el jugador eligiera la posición de inserción en el mazo:
\begin{enumerate}
    \item \textbf{Selector numérico}: Un control en el modal de acciones para indicar el índice exacto. Se descartó por añadir una carga visual excesiva y requerir validaciones adicionales de rango en el cliente.
    \item \textbf{Selección basada en UID}: Utilizar el diccionario de acciones pendientes para esperar el identificador de una carta ya presente en el mazo, permitiendo insertar la nueva carta justo antes de la seleccionada.
\end{enumerate}
Aunque se avanzó en la definición lógica de la segunda opción, se decidió finalmente delegar la implementación técnica completa y la creación del componente visual al otro integrante del equipo para garantizar la coherencia con el diseño de la interfaz del mazo.

\section{Detalles de Implementación}
El componente \textit{React} se programa para capturar el DTO efímero y congelar el entorno de la partida. Su estado visual se coordina estrictamente con la respuesta asíncrona del servidor, asegurando que los modales se desmonten limpiamente y reanuden el cauce natural de turnos tras la alteración exitosa en la base de datos de memoria.

La implementación técnica principal consiste en el desarrollo de la función auxiliar de cambio de turno en la capa de servicios. Esta función sustituye al método básico anterior, integrando una comprobación de rol y la llamada a \texttt{drawCardFromDeck} específicamente para el Rey.

\begin{lstlisting}[language=Java,basicstyle=\scriptsize\ttfamily,caption={Logica de cambio de turno con robo automático para el bando real},label={cod:king_draw_logic},breaklines=true]
const changeTurnWithDraw = async (gameState) => {
    if (gameState.turn === 'king') {
        drawCardFromDeck(gameState, 'king');
    }
    gameState.turn = (gameState.turn === 'king') ? 'peasant' : 'king';
    return gameState;
};
\end{lstlisting}

En cuanto a la infiltración, se prepararon los esqueletos de los diccionarios en el servidor (\texttt{peasantPendingActions}) para soportar el estado \texttt{INFILTRATE}, dejando la base preparada para que el compañero finalizara la integración con el componente gráfico avanzado.

\section{Seguimiento, Control y Replanificación}
Las pruebas del bucle final de juego levantan una alerta crítica: la invocación a \texttt{setupNewEra} corrompía la asignación de cartas en memoria al cambiar de fase. La desviación exige replanificar los ciclos del \textit{GameService}, reescribiendo desde cero las funciones de reasignación y limpieza de manos para asegurar que the nueva Era comience de forma determinista. Las tareas de automatización del flujo del Rey se completaron con éxito, eliminando errores de omisión de robo manual. Al transferir la responsabilidad de la mecánica de infiltración al otro desarrollador, se pudo optimizar el tiempo de este ciclo para centrarse en la robustez de las transiciones de estado, asegurando que el motor de juego no presentara inconsistencias al guardar los cambios en Redis y blindando el alcance ante posibles modificaciones externas del manual físico (\textbf{RP-01}). Reescribir las funciones de reasignación supuso una exigencia técnica que fue absorbida operativamente para no impactar las entregas. El esfuerzo garantizó el cierre definitivo de la tarea 4.2 (Lógica asimétrica de cartas y transiciones) de la EDT.


\chapter{Iteración 14: Pulido de Reglas, Corrección de Errores y Refactorización Estructural}

\section{Definición de Características}
El propósito central de este ciclo consiste en realizar un pulido exhaustivo de las reglas lógicas del motor de juego y en optimizar la estructura interna de la interfaz de usuario. Se busca perfeccionar los algoritmos de cálculo de las mecánicas principales, asegurando que la resolución de los enfrentamientos y los eventos de victoria se ejecuten con total precisión. Asimismo, se pretende llevar a cabo una reestructuración profunda de la arquitectura del \textit{frontend} para mejorar su escalabilidad y mantenibilidad, descomponiendo los elementos complejos en componentes más sencillos y reutilizables. Esta etapa busca elevar los estándares de calidad del software, corrigiendo anomalías funcionales y estableciendo un sistema de comunicación de errores más robusto que mejore la experiencia y la fiabilidad del sistema.

\section{Diseño de la Solución}
El diseño para refactorizar la excesiva complejidad del componente principal (\texttt{Game.tsx}) se fundamenta en tres pilares: extraer la lógica transaccional a \textit{Custom Hooks}, descomponer el árbol JSX en componentes atómicos (\texttt{PlayerArea.tsx}, \texttt{RivalArea.tsx}), y encapsular todas las interfaces de interrupción (\textit{Modals}) para mitigar re-renderizados.

Se rediseña el formateador de salida de estado (\texttt{saveAndFormatGameState()}) para que actúe como un interceptor pasivo. Cuando esta función empaca el estado para guardarlo y enviarlo, debe analizar de oficio si existen detonantes automáticos (como el Rey robando el \texttt{DECOY} infiltrado), forzando una pausa en el juego y un cambio de turno reactivo para que el oponente ejecute el efecto.

\section{Detalles de Implementación}
A nivel de reglas, se programan las correcciones matemáticas en los cálculos de letalidad. En la capa de \textit{frontend}, la refactorización se materializa separando el estado global de suscripciones a \textit{WebSockets} hacia \texttt{useGameActions.ts} y \texttt{useGameData.ts}. Adicionalmente, se codifica un nuevo sistema de notificaciones (\texttt{ErrorToast.tsx}) para interceptar y renderizar los errores enviados por el servidor.

Se corrigen validaciones en el \textit{frontend}, como impedir multiselecciones accidentales en cartas como \texttt{EXECUTOR}, prohibir seleccionar la propia carta ejecutada (error de \texttt{RAT}), arreglar opacidades gráficas (\textit{alpha channels}) en cartas no reveladas e impedir que los botones de menú se rendericen durante el bloqueo de una Acción Pendiente. En paralelo, la lógica del \texttt{DECOY} se empotra en el formateador del estado, ejecutándose silenciosamente tras cada robo del Rey.

\section{Seguimiento, Control y Replanificación}
Durante las pruebas de estrés de este ciclo, el control de errores caza un fallo 500 no controlado en el cierre de sesiones. Este defecto crítico se soluciona encapsulando el cierre de la partida en una gestión robusta de excepciones. Esta acción correctiva deja una base limpia y altamente estable. Con las reparaciones en las tarjetas de lógica compleja, el motor de juego probó ser altamente estable frente a interacciones de usuario erróneas. La encapsulación de estas excepciones se resolvió asumiendo una carga de trabajo adicional, manteniendo el proyecto en fecha y avanzando en la estabilización contemplada en el Paquete 4 de la EDT.


\chapter{Iteración 15: Notificaciones Globales, Reconexión Activa y Alertas de Partida}

\section{Definición de Características}
El objetivo principal de este ciclo consiste en desarrollar un sistema integral de notificaciones globales que mejore la comunicación de eventos críticos hacia el usuario. Se busca implementar mecanismos reactivos capaces de interceptar sucesos transversales, como la finalización de una partida o el abandono inesperado de un oponente, proporcionando una respuesta visual inmediata que guíe al jugador en su navegación. Asimismo, se pretende integrar una capa de control que coordine el flujo de información entre los diferentes módulos de la aplicación, asegurando que las notificaciones se presenten de forma coherente con el estado actual de la sesión. Esta etapa busca reforzar el ecosistema de comunicación del proyecto, garantizando que el usuario reciba un \textit{feedback} claro y preciso ante cualquier cambio científico en el entorno de juego.

Por otro lado, se establece como objetivo la gestión de inestabilidad de red y sistema de reconexión a la partida. Este consiste en la implementación de mecanismos de control en el servidor para gestionar de forma integral las deficiencias de conectividad de los usuarios. Esto abarca la detección y manejo de salidas abruptas y tiempos de espera excedidos (\textit{timeouts}), así como el desarrollo de una lógica de almacenamiento transitorio que facilita el retorno seguro e inmediato del jugador a la sala de espera o a la partida en curso tras una caída de red. El propósito de esta tarea técnica es proteger la integridad de la sesión, evitando que el flujo del juego quede bloqueado indefinidamente por la inactividad de un participante.

\section{Diseño de la Solución}
Dando cumplimiento al requisito de usabilidad, el diseño propone la creación del componente \texttt{AnnouncementModal.tsx}. La arquitectura planifica inyectarlo en la capa superior de la aplicación para interceptar los estados tanto en \texttt{Game.tsx} como en \texttt{LobbyRoom.tsx}, manteniéndolo oculto por defecto.

Se diseña un doble sistema de protección y reconexión activa para dar respuesta definitiva al riesgo crítico de pérdida de estado por desconexión abrupta (\textbf{RT-02}). Si un usuario se desconecta de una partida, se dispara un cronómetro de gracia de 2 minutos visible para el oponente. Si el tiempo expira, la victoria se adjudica automáticamente. Para la sala de espera (\textit{Lobby}), se diseña un \textit{hook} de validación que detecta si el usuario ya pertenece a una sesión activa y renderiza un panel de "Reconexión Rápida".

\section{Detalles de Implementación}
La integración se programa cambiando el estado global en los \textit{hooks} personalizados para escuchar de forma activa las interrupciones de \textit{Socket.io}. Al emitirse una notificación de victoria o abandono, el motor reactivo de estado dispara el modal, bloqueando la interfaz subyacente para forzar la navegación del usuario hacia una zona segura (como el menú principal) e invalidando las interacciones residuales.

Se genera un \texttt{useEffect()} anclado al objeto de sesión del usuario en la pantalla de inicio y listado. Si el servidor confirma la vinculación a una \texttt{LobbyId}, se renderizan componentes flotantes de acceso directo en lugar de procesar los botones estándar de creación. Los temporizadores del servidor se acoplaron a la desconexión del \textit{socket}, mutando el estado de victoria (\texttt{winner: id}) transcurridos los 120000 milisegundos de inactividad.

\section{Seguimiento, Control y Replanificación}
El seguimiento de la implementación reporta conflictos de renderizado: el modal retenía información obsoleta tras salidas abruptas de los \textit{Lobbies}. La solución requiere replanificar el ciclo de vida de desmontaje en \textit{React}, añadiendo lógicas estrictas para purgar las suscripciones a \textit{WebSockets} antes de destruir la vista. Tras esto, las notificaciones funcionan de forma impecable. Estas adiciones eliminaron los bloqueos severos donde los usuarios perdían el control de su sesión tras un refresco de la página. El flujo de contingencia cumple exitosamente con los estándares de usabilidad modernos. Tras purgar las suscripciones y absorber el impacto temporal mediante compresión de calendario, se certifica el cierre funcional de la tarea 4.5 (Construcción del tablero interactivo y animaciones de estado) según la estructura de la EDT.


\chapter{Iteración 16: Automatización de Pruebas, Cobertura de Código y Despliegue Final}

\section{Definición de Características}
El propósito fundamental de este ciclo consiste en garantizar la calidad técnica y la fiabilidad del ecosistema de \textit{King and Peasant} mediante procesos de validación automatizados. Se busca establecer una infraestructura de pruebas que permita blindar la lógica operativa del motor de juego y de los componentes de la interfaz de usuario, asegurando que el sistema responda de forma correcta ante diversos escenarios de interacción. Asimismo, se pretende someter al \textit{software} a una evaluación continua que identifique posibles fragilidades en la gestión de eventos asíncronos y en la integridad de los datos, permitiendo una corrección proactiva de anomalías estructurales.

Por otro lado, se busca preparar el entorno para su despliegue definitivo. Concretamente consiste en la adaptación y ejecución de \textit{scripts} automáticos encargados de inyectar los datos iniciales e indispensables (tales como la información base de las cartas o configuraciones predeterminadas) en la base de datos del entorno de despliegue real. Esta etapa busca certificar la estabilidad de la plataforma, consolidando un producto robusto y verificado en todas sus capas funcionales.

\section{Diseño de la Solución}
Se diseña una \textit{suite} de pruebas utilizando el motor \textit{Jest}. La arquitectura de validación simula servicios de usuarios y el \textit{Game Service} operando contra bases de datos aisladas en memoria. En \textit{frontend}, el diseño abarca pruebas de componentes mediante DOM virtual para verificar los flujos de inicio de sesión y gestión de perfil.

Se definió como meta de calidad técnica alcanzar al menos un 80\% de cobertura de código en los archivos de gestión de \textit{Lobbies}. Por otro lado, se diseñó una regla de validación adicional en la mecánica del Rey para imposibilitar físicamente que supere el límite reglamentario de 3 guardias desplegados.

\section{Detalles de Implementación}
Se redactan y ejecutan los casos de prueba mediante Jest y Vitest, encontrando agujeros de seguridad (por ejemplo, permitir que un usuario se uniera a una sala estando ya en otra, lo cual se bloqueó añadiendo un \textit{middleware} validador en la API). Se añadió una función comprobatoria en \texttt{placeCardInTown} que interrumpe la inserción de guardias si el contador supera 3. 

También se logra exponer fragilidades en lógicas de cascada (como \textit{Revolt}). La refactorización exige aislar las mutaciones de poder en buffers temporales. Se codifica una resolución en cascada estricta, asegurando que los eventos masivos no sobrescriban referencias de memoria. Esta aproximación garantiza que, si el sistema asíncrono detecta un fallo, la operación atómica aborte preservando la partida intacta.

Finalmente, se reescribe el \texttt{package.json} en la rutina de \textit{build} de Render para invocar \texttt{prisma db seed}, asegurando la carga obligatoria del catálogo de cartas oficial en el entorno en la nube.

\section{Seguimiento, Control y Replanificación}
Durante el monitoreo de la \textit{suite}, el sistema reporta discrepancias graves en \textit{Redis}: mecánicas letales (\textit{Strike}) no aplican el daño correctamente. El equipo pausa la creación de nuevos \textit{tests} y replanificó el esfuerzo para reescribir parte de la lógica en los \textit{pendingActions} defectuosos. Tras iterar sobre el código, la tasa de éxito de las pruebas alcanza el 100\%, validando la estabilidad definitiva del software. Los tests cumplieron su propósito revelando deficiencias de borde en las rutas HTTP. Con la base de datos de producción poblándose automáticamente y la cobertura de pruebas establecida, la robustez general de la aplicación quedó certificada y lista para su defensa y evaluación. La reescritura de los 	extit{tests} defectuosos requirió una elevada compensación horaria en la recta final. A nivel organizativo, se certifica el cierre de la tarea 4.6 (Baterías de pruebas y estabilización) y se valida definitivamente la tarea 1.3 (Flujos de Integración Continua) de la EDT.


\chapter{Iteración 17: Redacción Técnica y Consolidación Documental}

\section{Definición de Características}
El objetivo primordial de este ciclo final consiste en intensificar la producción de la memoria técnica del proyecto, liberando ancho de banda y presión administrativa para afrontar la entrega definitiva. Se busca integrar y refinar los registros técnicos y manuales generados de forma incremental durante las iteraciones previas, asegurando que la memoria técnica presente una estructura cohesiva y profesional. Asimismo, se pretende llevar a cabo una revisión exhaustiva de la trazabilidad entre los requisitos iniciales y la implementación final, verificando que todos los artefactos visuales y diagramas de diseño estén correctamente alineados con el estado actual del proyecto para su entrega definitiva.

\section{Diseño de la Solución}
La solución metodológica plantea detener el desarrollo de nuevas características y establecer un flujo de redacción concurrente. El diseño documental exige justificar las decisiones arquitectónicas previas y estructurar todos los artefactos visuales (diagramas UML, tablas de análisis) generados a lo largo de los 16 \textit{sprints} anteriores.

\section{Detalles de Implementación}
La ejecución se centra en el trabajo sobre el proyecto \textit{LaTeX}. Se unifica el estilo narrativo del equipo y se resuelven errores tipográficos y de compilación (\textit{Overfull hboxes}). La redacción se aborda con revisiones cruzadas constantes frente al repositorio de código fuente para garantizar que la memoria sea un reflejo técnico exacto de la ingeniería de \textit{King and Peasant}.

\section{Seguimiento, Control y Replanificación}
El control de la iteración no reporta bloqueos de software, pero evidencia la carga cognitiva de transicionar del código a la redacción académica. La gestión del tiempo se ajusta mediante sesiones cortas de validación, logrando completar la estructura de la memoria sin desviaciones temporales. El ciclo cierra el cronograma de iteraciones con un éxito rotundo. El avance técnico alcanza la totalidad de sus objetivos, clausurando formalmente la tarea 1.4 (Elaboración de la memoria técnica) ilustrada en la EDT y dando por concluido el ciclo de desarrollo del proyecto.